% =============================================================================
% autoexam template — randomized, multi-version exam with key
% =============================================================================
%
% Build instructor copy (single version):
%   lualatex \def\Version{A}% =============================================================================
% autoexam template — randomized, multi-version exam with key
% =============================================================================
%
% Build instructor copy (single version):
%   lualatex \def\Version{A}% =============================================================================
% autoexam template — randomized, multi-version exam with key
% =============================================================================
%
% Build instructor copy (single version):
%   lualatex \def\Version{A}% =============================================================================
% autoexam template — randomized, multi-version exam with key
% =============================================================================
%
% Build instructor copy (single version):
%   lualatex \def\Version{A}\input{autoexam-template.tex}
%
% Build full collated multi-version PDF (no \Version):
%   lualatex autoexam-template.tex
%
% Build answer key:
%   lualatex \def\ShowKey{}\input{autoexam-template.tex}
%
% Build with rubric overlays:
%   lualatex \def\ShowRubric{}\input{autoexam-template.tex}

\documentclass[
	exam-number     = 5,
	exam-date       = {May 2, 2026},
	exam-postscript = {Good luck!},
]{autoexam}

% -- Versions ------------------------------------------------------------------
% Declare the per-exam letter labels. The Python builder (if used) parses
% this line via regex to discover versions.
\versions{A, B, C}

% -- Problem bank --------------------------------------------------------------
% \loadbank pulls in problem definitions; alternatively, you can put
% \begin{problem}{id} ... \end{problem} blocks here (in the preamble or body).
\loadbank{bank.tex}

\begin{document}

% -- Cover page ----------------------------------------------------------------
% \maketitle emits the exam cover: the title (Exam <number>), course and date,
% the version circle (shown when \versions is declared), the boxed instructions
% (from autoexam-instructions.tex unless overridden), and the student name line.
% \begin{problems} issues a \clearpage after it, so Q1 starts on a fresh page.
% Add \blankpage here for a pre-problem scratch-work leaf.
\maketitle

% -- Body ----------------------------------------------------------------------
% Problems are grouped via \begin{problems} ... \end{problems}, which
% performs the smart two-column layout based on per-problem hints.
\begin{problems}

	% Inside \begin{problems}, use \problem{...} -- it emits the \question
	% wrapper that a problem's \begin{parts} needs. (\getproblem emits only
	% the bare problem body and is meant to follow an explicit \question,
	% e.g. in the quiz template.) The argument is a bank id or a key=value
	% filter; an optional [p1,p2,...] distributes points to \ppart calls.
	\problem{topic=quad}
	\problem{topic=ratineq}
	\problem{topic=graph}

	% An inline (one-off) problem can be placed right here:
	%
	% \begin{problem}{linear_eq}[topic=algebra, diff=easy]
	%     Solve $2x + 5 = 11$ for $x$.
	%     \solution
	%     $x = 3$.
	% \end{problem}

\end{problems}

% Optional: a final scoring grid for the front desk
% \scorepage[20]

\end{document}

%
% Build full collated multi-version PDF (no \Version):
%   lualatex autoexam-template.tex
%
% Build answer key:
%   lualatex \def\ShowKey{}% =============================================================================
% autoexam template — randomized, multi-version exam with key
% =============================================================================
%
% Build instructor copy (single version):
%   lualatex \def\Version{A}\input{autoexam-template.tex}
%
% Build full collated multi-version PDF (no \Version):
%   lualatex autoexam-template.tex
%
% Build answer key:
%   lualatex \def\ShowKey{}\input{autoexam-template.tex}
%
% Build with rubric overlays:
%   lualatex \def\ShowRubric{}\input{autoexam-template.tex}

\documentclass[
	exam-number     = 5,
	exam-date       = {May 2, 2026},
	exam-postscript = {Good luck!},
]{autoexam}

% -- Versions ------------------------------------------------------------------
% Declare the per-exam letter labels. The Python builder (if used) parses
% this line via regex to discover versions.
\versions{A, B, C}

% -- Problem bank --------------------------------------------------------------
% \loadbank pulls in problem definitions; alternatively, you can put
% \begin{problem}{id} ... \end{problem} blocks here (in the preamble or body).
\loadbank{bank.tex}

\begin{document}

% -- Cover page ----------------------------------------------------------------
% \maketitle emits the exam cover: the title (Exam <number>), course and date,
% the version circle (shown when \versions is declared), the boxed instructions
% (from autoexam-instructions.tex unless overridden), and the student name line.
% \begin{problems} issues a \clearpage after it, so Q1 starts on a fresh page.
% Add \blankpage here for a pre-problem scratch-work leaf.
\maketitle

% -- Body ----------------------------------------------------------------------
% Problems are grouped via \begin{problems} ... \end{problems}, which
% performs the smart two-column layout based on per-problem hints.
\begin{problems}

	% Inside \begin{problems}, use \problem{...} -- it emits the \question
	% wrapper that a problem's \begin{parts} needs. (\getproblem emits only
	% the bare problem body and is meant to follow an explicit \question,
	% e.g. in the quiz template.) The argument is a bank id or a key=value
	% filter; an optional [p1,p2,...] distributes points to \ppart calls.
	\problem{topic=quad}
	\problem{topic=ratineq}
	\problem{topic=graph}

	% An inline (one-off) problem can be placed right here:
	%
	% \begin{problem}{linear_eq}[topic=algebra, diff=easy]
	%     Solve $2x + 5 = 11$ for $x$.
	%     \solution
	%     $x = 3$.
	% \end{problem}

\end{problems}

% Optional: a final scoring grid for the front desk
% \scorepage[20]

\end{document}

%
% Build with rubric overlays:
%   lualatex \def\ShowRubric{}% =============================================================================
% autoexam template — randomized, multi-version exam with key
% =============================================================================
%
% Build instructor copy (single version):
%   lualatex \def\Version{A}\input{autoexam-template.tex}
%
% Build full collated multi-version PDF (no \Version):
%   lualatex autoexam-template.tex
%
% Build answer key:
%   lualatex \def\ShowKey{}\input{autoexam-template.tex}
%
% Build with rubric overlays:
%   lualatex \def\ShowRubric{}\input{autoexam-template.tex}

\documentclass[
	exam-number     = 5,
	exam-date       = {May 2, 2026},
	exam-postscript = {Good luck!},
]{autoexam}

% -- Versions ------------------------------------------------------------------
% Declare the per-exam letter labels. The Python builder (if used) parses
% this line via regex to discover versions.
\versions{A, B, C}

% -- Problem bank --------------------------------------------------------------
% \loadbank pulls in problem definitions; alternatively, you can put
% \begin{problem}{id} ... \end{problem} blocks here (in the preamble or body).
\loadbank{bank.tex}

\begin{document}

% -- Cover page ----------------------------------------------------------------
% \maketitle emits the exam cover: the title (Exam <number>), course and date,
% the version circle (shown when \versions is declared), the boxed instructions
% (from autoexam-instructions.tex unless overridden), and the student name line.
% \begin{problems} issues a \clearpage after it, so Q1 starts on a fresh page.
% Add \blankpage here for a pre-problem scratch-work leaf.
\maketitle

% -- Body ----------------------------------------------------------------------
% Problems are grouped via \begin{problems} ... \end{problems}, which
% performs the smart two-column layout based on per-problem hints.
\begin{problems}

	% Inside \begin{problems}, use \problem{...} -- it emits the \question
	% wrapper that a problem's \begin{parts} needs. (\getproblem emits only
	% the bare problem body and is meant to follow an explicit \question,
	% e.g. in the quiz template.) The argument is a bank id or a key=value
	% filter; an optional [p1,p2,...] distributes points to \ppart calls.
	\problem{topic=quad}
	\problem{topic=ratineq}
	\problem{topic=graph}

	% An inline (one-off) problem can be placed right here:
	%
	% \begin{problem}{linear_eq}[topic=algebra, diff=easy]
	%     Solve $2x + 5 = 11$ for $x$.
	%     \solution
	%     $x = 3$.
	% \end{problem}

\end{problems}

% Optional: a final scoring grid for the front desk
% \scorepage[20]

\end{document}


\documentclass[
	exam-number     = 5,
	exam-date       = {May 2, 2026},
	exam-postscript = {Good luck!},
]{autoexam}

% -- Versions ------------------------------------------------------------------
% Declare the per-exam letter labels. The Python builder (if used) parses
% this line via regex to discover versions.
\versions{A, B, C}

% -- Problem bank --------------------------------------------------------------
% \loadbank pulls in problem definitions; alternatively, you can put
% \begin{problem}{id} ... \end{problem} blocks here (in the preamble or body).
\loadbank{bank.tex}

\begin{document}

% -- Cover page ----------------------------------------------------------------
% \maketitle emits the exam cover: the title (Exam <number>), course and date,
% the version circle (shown when \versions is declared), the boxed instructions
% (from autoexam-instructions.tex unless overridden), and the student name line.
% \begin{problems} issues a \clearpage after it, so Q1 starts on a fresh page.
% Add \blankpage here for a pre-problem scratch-work leaf.
\maketitle

% -- Body ----------------------------------------------------------------------
% Problems are grouped via \begin{problems} ... \end{problems}, which
% performs the smart two-column layout based on per-problem hints.
\begin{problems}

	% Inside \begin{problems}, use \problem{...} -- it emits the \question
	% wrapper that a problem's \begin{parts} needs. (\getproblem emits only
	% the bare problem body and is meant to follow an explicit \question,
	% e.g. in the quiz template.) The argument is a bank id or a key=value
	% filter; an optional [p1,p2,...] distributes points to \ppart calls.
	\problem{topic=quad}
	\problem{topic=ratineq}
	\problem{topic=graph}

	% An inline (one-off) problem can be placed right here:
	%
	% \begin{problem}{linear_eq}[topic=algebra, diff=easy]
	%     Solve $2x + 5 = 11$ for $x$.
	%     \solution
	%     $x = 3$.
	% \end{problem}

\end{problems}

% Optional: a final scoring grid for the front desk
% \scorepage[20]

\end{document}

%
% Build full collated multi-version PDF (no \Version):
%   lualatex autoexam-template.tex
%
% Build answer key:
%   lualatex \def\ShowKey{}% =============================================================================
% autoexam template — randomized, multi-version exam with key
% =============================================================================
%
% Build instructor copy (single version):
%   lualatex \def\Version{A}% =============================================================================
% autoexam template — randomized, multi-version exam with key
% =============================================================================
%
% Build instructor copy (single version):
%   lualatex \def\Version{A}\input{autoexam-template.tex}
%
% Build full collated multi-version PDF (no \Version):
%   lualatex autoexam-template.tex
%
% Build answer key:
%   lualatex \def\ShowKey{}\input{autoexam-template.tex}
%
% Build with rubric overlays:
%   lualatex \def\ShowRubric{}\input{autoexam-template.tex}

\documentclass[
	exam-number     = 5,
	exam-date       = {May 2, 2026},
	exam-postscript = {Good luck!},
]{autoexam}

% -- Versions ------------------------------------------------------------------
% Declare the per-exam letter labels. The Python builder (if used) parses
% this line via regex to discover versions.
\versions{A, B, C}

% -- Problem bank --------------------------------------------------------------
% \loadbank pulls in problem definitions; alternatively, you can put
% \begin{problem}{id} ... \end{problem} blocks here (in the preamble or body).
\loadbank{bank.tex}

\begin{document}

% -- Cover page ----------------------------------------------------------------
% \maketitle emits the exam cover: the title (Exam <number>), course and date,
% the version circle (shown when \versions is declared), the boxed instructions
% (from autoexam-instructions.tex unless overridden), and the student name line.
% \begin{problems} issues a \clearpage after it, so Q1 starts on a fresh page.
% Add \blankpage here for a pre-problem scratch-work leaf.
\maketitle

% -- Body ----------------------------------------------------------------------
% Problems are grouped via \begin{problems} ... \end{problems}, which
% performs the smart two-column layout based on per-problem hints.
\begin{problems}

	% Inside \begin{problems}, use \problem{...} -- it emits the \question
	% wrapper that a problem's \begin{parts} needs. (\getproblem emits only
	% the bare problem body and is meant to follow an explicit \question,
	% e.g. in the quiz template.) The argument is a bank id or a key=value
	% filter; an optional [p1,p2,...] distributes points to \ppart calls.
	\problem{topic=quad}
	\problem{topic=ratineq}
	\problem{topic=graph}

	% An inline (one-off) problem can be placed right here:
	%
	% \begin{problem}{linear_eq}[topic=algebra, diff=easy]
	%     Solve $2x + 5 = 11$ for $x$.
	%     \solution
	%     $x = 3$.
	% \end{problem}

\end{problems}

% Optional: a final scoring grid for the front desk
% \scorepage[20]

\end{document}

%
% Build full collated multi-version PDF (no \Version):
%   lualatex autoexam-template.tex
%
% Build answer key:
%   lualatex \def\ShowKey{}% =============================================================================
% autoexam template — randomized, multi-version exam with key
% =============================================================================
%
% Build instructor copy (single version):
%   lualatex \def\Version{A}\input{autoexam-template.tex}
%
% Build full collated multi-version PDF (no \Version):
%   lualatex autoexam-template.tex
%
% Build answer key:
%   lualatex \def\ShowKey{}\input{autoexam-template.tex}
%
% Build with rubric overlays:
%   lualatex \def\ShowRubric{}\input{autoexam-template.tex}

\documentclass[
	exam-number     = 5,
	exam-date       = {May 2, 2026},
	exam-postscript = {Good luck!},
]{autoexam}

% -- Versions ------------------------------------------------------------------
% Declare the per-exam letter labels. The Python builder (if used) parses
% this line via regex to discover versions.
\versions{A, B, C}

% -- Problem bank --------------------------------------------------------------
% \loadbank pulls in problem definitions; alternatively, you can put
% \begin{problem}{id} ... \end{problem} blocks here (in the preamble or body).
\loadbank{bank.tex}

\begin{document}

% -- Cover page ----------------------------------------------------------------
% \maketitle emits the exam cover: the title (Exam <number>), course and date,
% the version circle (shown when \versions is declared), the boxed instructions
% (from autoexam-instructions.tex unless overridden), and the student name line.
% \begin{problems} issues a \clearpage after it, so Q1 starts on a fresh page.
% Add \blankpage here for a pre-problem scratch-work leaf.
\maketitle

% -- Body ----------------------------------------------------------------------
% Problems are grouped via \begin{problems} ... \end{problems}, which
% performs the smart two-column layout based on per-problem hints.
\begin{problems}

	% Inside \begin{problems}, use \problem{...} -- it emits the \question
	% wrapper that a problem's \begin{parts} needs. (\getproblem emits only
	% the bare problem body and is meant to follow an explicit \question,
	% e.g. in the quiz template.) The argument is a bank id or a key=value
	% filter; an optional [p1,p2,...] distributes points to \ppart calls.
	\problem{topic=quad}
	\problem{topic=ratineq}
	\problem{topic=graph}

	% An inline (one-off) problem can be placed right here:
	%
	% \begin{problem}{linear_eq}[topic=algebra, diff=easy]
	%     Solve $2x + 5 = 11$ for $x$.
	%     \solution
	%     $x = 3$.
	% \end{problem}

\end{problems}

% Optional: a final scoring grid for the front desk
% \scorepage[20]

\end{document}

%
% Build with rubric overlays:
%   lualatex \def\ShowRubric{}% =============================================================================
% autoexam template — randomized, multi-version exam with key
% =============================================================================
%
% Build instructor copy (single version):
%   lualatex \def\Version{A}\input{autoexam-template.tex}
%
% Build full collated multi-version PDF (no \Version):
%   lualatex autoexam-template.tex
%
% Build answer key:
%   lualatex \def\ShowKey{}\input{autoexam-template.tex}
%
% Build with rubric overlays:
%   lualatex \def\ShowRubric{}\input{autoexam-template.tex}

\documentclass[
	exam-number     = 5,
	exam-date       = {May 2, 2026},
	exam-postscript = {Good luck!},
]{autoexam}

% -- Versions ------------------------------------------------------------------
% Declare the per-exam letter labels. The Python builder (if used) parses
% this line via regex to discover versions.
\versions{A, B, C}

% -- Problem bank --------------------------------------------------------------
% \loadbank pulls in problem definitions; alternatively, you can put
% \begin{problem}{id} ... \end{problem} blocks here (in the preamble or body).
\loadbank{bank.tex}

\begin{document}

% -- Cover page ----------------------------------------------------------------
% \maketitle emits the exam cover: the title (Exam <number>), course and date,
% the version circle (shown when \versions is declared), the boxed instructions
% (from autoexam-instructions.tex unless overridden), and the student name line.
% \begin{problems} issues a \clearpage after it, so Q1 starts on a fresh page.
% Add \blankpage here for a pre-problem scratch-work leaf.
\maketitle

% -- Body ----------------------------------------------------------------------
% Problems are grouped via \begin{problems} ... \end{problems}, which
% performs the smart two-column layout based on per-problem hints.
\begin{problems}

	% Inside \begin{problems}, use \problem{...} -- it emits the \question
	% wrapper that a problem's \begin{parts} needs. (\getproblem emits only
	% the bare problem body and is meant to follow an explicit \question,
	% e.g. in the quiz template.) The argument is a bank id or a key=value
	% filter; an optional [p1,p2,...] distributes points to \ppart calls.
	\problem{topic=quad}
	\problem{topic=ratineq}
	\problem{topic=graph}

	% An inline (one-off) problem can be placed right here:
	%
	% \begin{problem}{linear_eq}[topic=algebra, diff=easy]
	%     Solve $2x + 5 = 11$ for $x$.
	%     \solution
	%     $x = 3$.
	% \end{problem}

\end{problems}

% Optional: a final scoring grid for the front desk
% \scorepage[20]

\end{document}


\documentclass[
	exam-number     = 5,
	exam-date       = {May 2, 2026},
	exam-postscript = {Good luck!},
]{autoexam}

% -- Versions ------------------------------------------------------------------
% Declare the per-exam letter labels. The Python builder (if used) parses
% this line via regex to discover versions.
\versions{A, B, C}

% -- Problem bank --------------------------------------------------------------
% \loadbank pulls in problem definitions; alternatively, you can put
% \begin{problem}{id} ... \end{problem} blocks here (in the preamble or body).
\loadbank{bank.tex}

\begin{document}

% -- Cover page ----------------------------------------------------------------
% \maketitle emits the exam cover: the title (Exam <number>), course and date,
% the version circle (shown when \versions is declared), the boxed instructions
% (from autoexam-instructions.tex unless overridden), and the student name line.
% \begin{problems} issues a \clearpage after it, so Q1 starts on a fresh page.
% Add \blankpage here for a pre-problem scratch-work leaf.
\maketitle

% -- Body ----------------------------------------------------------------------
% Problems are grouped via \begin{problems} ... \end{problems}, which
% performs the smart two-column layout based on per-problem hints.
\begin{problems}

	% Inside \begin{problems}, use \problem{...} -- it emits the \question
	% wrapper that a problem's \begin{parts} needs. (\getproblem emits only
	% the bare problem body and is meant to follow an explicit \question,
	% e.g. in the quiz template.) The argument is a bank id or a key=value
	% filter; an optional [p1,p2,...] distributes points to \ppart calls.
	\problem{topic=quad}
	\problem{topic=ratineq}
	\problem{topic=graph}

	% An inline (one-off) problem can be placed right here:
	%
	% \begin{problem}{linear_eq}[topic=algebra, diff=easy]
	%     Solve $2x + 5 = 11$ for $x$.
	%     \solution
	%     $x = 3$.
	% \end{problem}

\end{problems}

% Optional: a final scoring grid for the front desk
% \scorepage[20]

\end{document}

%
% Build with rubric overlays:
%   lualatex \def\ShowRubric{}% =============================================================================
% autoexam template — randomized, multi-version exam with key
% =============================================================================
%
% Build instructor copy (single version):
%   lualatex \def\Version{A}% =============================================================================
% autoexam template — randomized, multi-version exam with key
% =============================================================================
%
% Build instructor copy (single version):
%   lualatex \def\Version{A}\input{autoexam-template.tex}
%
% Build full collated multi-version PDF (no \Version):
%   lualatex autoexam-template.tex
%
% Build answer key:
%   lualatex \def\ShowKey{}\input{autoexam-template.tex}
%
% Build with rubric overlays:
%   lualatex \def\ShowRubric{}\input{autoexam-template.tex}

\documentclass[
	exam-number     = 5,
	exam-date       = {May 2, 2026},
	exam-postscript = {Good luck!},
]{autoexam}

% -- Versions ------------------------------------------------------------------
% Declare the per-exam letter labels. The Python builder (if used) parses
% this line via regex to discover versions.
\versions{A, B, C}

% -- Problem bank --------------------------------------------------------------
% \loadbank pulls in problem definitions; alternatively, you can put
% \begin{problem}{id} ... \end{problem} blocks here (in the preamble or body).
\loadbank{bank.tex}

\begin{document}

% -- Cover page ----------------------------------------------------------------
% \maketitle emits the exam cover: the title (Exam <number>), course and date,
% the version circle (shown when \versions is declared), the boxed instructions
% (from autoexam-instructions.tex unless overridden), and the student name line.
% \begin{problems} issues a \clearpage after it, so Q1 starts on a fresh page.
% Add \blankpage here for a pre-problem scratch-work leaf.
\maketitle

% -- Body ----------------------------------------------------------------------
% Problems are grouped via \begin{problems} ... \end{problems}, which
% performs the smart two-column layout based on per-problem hints.
\begin{problems}

	% Inside \begin{problems}, use \problem{...} -- it emits the \question
	% wrapper that a problem's \begin{parts} needs. (\getproblem emits only
	% the bare problem body and is meant to follow an explicit \question,
	% e.g. in the quiz template.) The argument is a bank id or a key=value
	% filter; an optional [p1,p2,...] distributes points to \ppart calls.
	\problem{topic=quad}
	\problem{topic=ratineq}
	\problem{topic=graph}

	% An inline (one-off) problem can be placed right here:
	%
	% \begin{problem}{linear_eq}[topic=algebra, diff=easy]
	%     Solve $2x + 5 = 11$ for $x$.
	%     \solution
	%     $x = 3$.
	% \end{problem}

\end{problems}

% Optional: a final scoring grid for the front desk
% \scorepage[20]

\end{document}

%
% Build full collated multi-version PDF (no \Version):
%   lualatex autoexam-template.tex
%
% Build answer key:
%   lualatex \def\ShowKey{}% =============================================================================
% autoexam template — randomized, multi-version exam with key
% =============================================================================
%
% Build instructor copy (single version):
%   lualatex \def\Version{A}\input{autoexam-template.tex}
%
% Build full collated multi-version PDF (no \Version):
%   lualatex autoexam-template.tex
%
% Build answer key:
%   lualatex \def\ShowKey{}\input{autoexam-template.tex}
%
% Build with rubric overlays:
%   lualatex \def\ShowRubric{}\input{autoexam-template.tex}

\documentclass[
	exam-number     = 5,
	exam-date       = {May 2, 2026},
	exam-postscript = {Good luck!},
]{autoexam}

% -- Versions ------------------------------------------------------------------
% Declare the per-exam letter labels. The Python builder (if used) parses
% this line via regex to discover versions.
\versions{A, B, C}

% -- Problem bank --------------------------------------------------------------
% \loadbank pulls in problem definitions; alternatively, you can put
% \begin{problem}{id} ... \end{problem} blocks here (in the preamble or body).
\loadbank{bank.tex}

\begin{document}

% -- Cover page ----------------------------------------------------------------
% \maketitle emits the exam cover: the title (Exam <number>), course and date,
% the version circle (shown when \versions is declared), the boxed instructions
% (from autoexam-instructions.tex unless overridden), and the student name line.
% \begin{problems} issues a \clearpage after it, so Q1 starts on a fresh page.
% Add \blankpage here for a pre-problem scratch-work leaf.
\maketitle

% -- Body ----------------------------------------------------------------------
% Problems are grouped via \begin{problems} ... \end{problems}, which
% performs the smart two-column layout based on per-problem hints.
\begin{problems}

	% Inside \begin{problems}, use \problem{...} -- it emits the \question
	% wrapper that a problem's \begin{parts} needs. (\getproblem emits only
	% the bare problem body and is meant to follow an explicit \question,
	% e.g. in the quiz template.) The argument is a bank id or a key=value
	% filter; an optional [p1,p2,...] distributes points to \ppart calls.
	\problem{topic=quad}
	\problem{topic=ratineq}
	\problem{topic=graph}

	% An inline (one-off) problem can be placed right here:
	%
	% \begin{problem}{linear_eq}[topic=algebra, diff=easy]
	%     Solve $2x + 5 = 11$ for $x$.
	%     \solution
	%     $x = 3$.
	% \end{problem}

\end{problems}

% Optional: a final scoring grid for the front desk
% \scorepage[20]

\end{document}

%
% Build with rubric overlays:
%   lualatex \def\ShowRubric{}% =============================================================================
% autoexam template — randomized, multi-version exam with key
% =============================================================================
%
% Build instructor copy (single version):
%   lualatex \def\Version{A}\input{autoexam-template.tex}
%
% Build full collated multi-version PDF (no \Version):
%   lualatex autoexam-template.tex
%
% Build answer key:
%   lualatex \def\ShowKey{}\input{autoexam-template.tex}
%
% Build with rubric overlays:
%   lualatex \def\ShowRubric{}\input{autoexam-template.tex}

\documentclass[
	exam-number     = 5,
	exam-date       = {May 2, 2026},
	exam-postscript = {Good luck!},
]{autoexam}

% -- Versions ------------------------------------------------------------------
% Declare the per-exam letter labels. The Python builder (if used) parses
% this line via regex to discover versions.
\versions{A, B, C}

% -- Problem bank --------------------------------------------------------------
% \loadbank pulls in problem definitions; alternatively, you can put
% \begin{problem}{id} ... \end{problem} blocks here (in the preamble or body).
\loadbank{bank.tex}

\begin{document}

% -- Cover page ----------------------------------------------------------------
% \maketitle emits the exam cover: the title (Exam <number>), course and date,
% the version circle (shown when \versions is declared), the boxed instructions
% (from autoexam-instructions.tex unless overridden), and the student name line.
% \begin{problems} issues a \clearpage after it, so Q1 starts on a fresh page.
% Add \blankpage here for a pre-problem scratch-work leaf.
\maketitle

% -- Body ----------------------------------------------------------------------
% Problems are grouped via \begin{problems} ... \end{problems}, which
% performs the smart two-column layout based on per-problem hints.
\begin{problems}

	% Inside \begin{problems}, use \problem{...} -- it emits the \question
	% wrapper that a problem's \begin{parts} needs. (\getproblem emits only
	% the bare problem body and is meant to follow an explicit \question,
	% e.g. in the quiz template.) The argument is a bank id or a key=value
	% filter; an optional [p1,p2,...] distributes points to \ppart calls.
	\problem{topic=quad}
	\problem{topic=ratineq}
	\problem{topic=graph}

	% An inline (one-off) problem can be placed right here:
	%
	% \begin{problem}{linear_eq}[topic=algebra, diff=easy]
	%     Solve $2x + 5 = 11$ for $x$.
	%     \solution
	%     $x = 3$.
	% \end{problem}

\end{problems}

% Optional: a final scoring grid for the front desk
% \scorepage[20]

\end{document}


\documentclass[
	exam-number     = 5,
	exam-date       = {May 2, 2026},
	exam-postscript = {Good luck!},
]{autoexam}

% -- Versions ------------------------------------------------------------------
% Declare the per-exam letter labels. The Python builder (if used) parses
% this line via regex to discover versions.
\versions{A, B, C}

% -- Problem bank --------------------------------------------------------------
% \loadbank pulls in problem definitions; alternatively, you can put
% \begin{problem}{id} ... \end{problem} blocks here (in the preamble or body).
\loadbank{bank.tex}

\begin{document}

% -- Cover page ----------------------------------------------------------------
% \maketitle emits the exam cover: the title (Exam <number>), course and date,
% the version circle (shown when \versions is declared), the boxed instructions
% (from autoexam-instructions.tex unless overridden), and the student name line.
% \begin{problems} issues a \clearpage after it, so Q1 starts on a fresh page.
% Add \blankpage here for a pre-problem scratch-work leaf.
\maketitle

% -- Body ----------------------------------------------------------------------
% Problems are grouped via \begin{problems} ... \end{problems}, which
% performs the smart two-column layout based on per-problem hints.
\begin{problems}

	% Inside \begin{problems}, use \problem{...} -- it emits the \question
	% wrapper that a problem's \begin{parts} needs. (\getproblem emits only
	% the bare problem body and is meant to follow an explicit \question,
	% e.g. in the quiz template.) The argument is a bank id or a key=value
	% filter; an optional [p1,p2,...] distributes points to \ppart calls.
	\problem{topic=quad}
	\problem{topic=ratineq}
	\problem{topic=graph}

	% An inline (one-off) problem can be placed right here:
	%
	% \begin{problem}{linear_eq}[topic=algebra, diff=easy]
	%     Solve $2x + 5 = 11$ for $x$.
	%     \solution
	%     $x = 3$.
	% \end{problem}

\end{problems}

% Optional: a final scoring grid for the front desk
% \scorepage[20]

\end{document}


\documentclass[
	exam-number     = 5,
	exam-date       = {May 2, 2026},
	exam-postscript = {Good luck!},
]{autoexam}

% -- Versions ------------------------------------------------------------------
% Declare the per-exam letter labels. The Python builder (if used) parses
% this line via regex to discover versions.
\versions{A, B, C}

% -- Problem bank --------------------------------------------------------------
% \loadbank pulls in problem definitions; alternatively, you can put
% \begin{problem}{id} ... \end{problem} blocks here (in the preamble or body).
\loadbank{bank.tex}

\begin{document}

% -- Cover page ----------------------------------------------------------------
% \maketitle emits the exam cover: the title (Exam <number>), course and date,
% the version circle (shown when \versions is declared), the boxed instructions
% (from autoexam-instructions.tex unless overridden), and the student name line.
% \begin{problems} issues a \clearpage after it, so Q1 starts on a fresh page.
% Add \blankpage here for a pre-problem scratch-work leaf.
\maketitle

% -- Body ----------------------------------------------------------------------
% Problems are grouped via \begin{problems} ... \end{problems}, which
% performs the smart two-column layout based on per-problem hints.
\begin{problems}

	% Inside \begin{problems}, use \problem{...} -- it emits the \question
	% wrapper that a problem's \begin{parts} needs. (\getproblem emits only
	% the bare problem body and is meant to follow an explicit \question,
	% e.g. in the quiz template.) The argument is a bank id or a key=value
	% filter; an optional [p1,p2,...] distributes points to \ppart calls.
	\problem{topic=quad}
	\problem{topic=ratineq}
	\problem{topic=graph}

	% An inline (one-off) problem can be placed right here:
	%
	% \begin{problem}{linear_eq}[topic=algebra, diff=easy]
	%     Solve $2x + 5 = 11$ for $x$.
	%     \solution
	%     $x = 3$.
	% \end{problem}

\end{problems}

% Optional: a final scoring grid for the front desk
% \scorepage[20]

\end{document}

%
% Build full collated multi-version PDF (no \Version):
%   lualatex autoexam-template.tex
%
% Build answer key:
%   lualatex \def\ShowKey{}% =============================================================================
% autoexam template — randomized, multi-version exam with key
% =============================================================================
%
% Build instructor copy (single version):
%   lualatex \def\Version{A}% =============================================================================
% autoexam template — randomized, multi-version exam with key
% =============================================================================
%
% Build instructor copy (single version):
%   lualatex \def\Version{A}% =============================================================================
% autoexam template — randomized, multi-version exam with key
% =============================================================================
%
% Build instructor copy (single version):
%   lualatex \def\Version{A}\input{autoexam-template.tex}
%
% Build full collated multi-version PDF (no \Version):
%   lualatex autoexam-template.tex
%
% Build answer key:
%   lualatex \def\ShowKey{}\input{autoexam-template.tex}
%
% Build with rubric overlays:
%   lualatex \def\ShowRubric{}\input{autoexam-template.tex}

\documentclass[
	exam-number     = 5,
	exam-date       = {May 2, 2026},
	exam-postscript = {Good luck!},
]{autoexam}

% -- Versions ------------------------------------------------------------------
% Declare the per-exam letter labels. The Python builder (if used) parses
% this line via regex to discover versions.
\versions{A, B, C}

% -- Problem bank --------------------------------------------------------------
% \loadbank pulls in problem definitions; alternatively, you can put
% \begin{problem}{id} ... \end{problem} blocks here (in the preamble or body).
\loadbank{bank.tex}

\begin{document}

% -- Cover page ----------------------------------------------------------------
% \maketitle emits the exam cover: the title (Exam <number>), course and date,
% the version circle (shown when \versions is declared), the boxed instructions
% (from autoexam-instructions.tex unless overridden), and the student name line.
% \begin{problems} issues a \clearpage after it, so Q1 starts on a fresh page.
% Add \blankpage here for a pre-problem scratch-work leaf.
\maketitle

% -- Body ----------------------------------------------------------------------
% Problems are grouped via \begin{problems} ... \end{problems}, which
% performs the smart two-column layout based on per-problem hints.
\begin{problems}

	% Inside \begin{problems}, use \problem{...} -- it emits the \question
	% wrapper that a problem's \begin{parts} needs. (\getproblem emits only
	% the bare problem body and is meant to follow an explicit \question,
	% e.g. in the quiz template.) The argument is a bank id or a key=value
	% filter; an optional [p1,p2,...] distributes points to \ppart calls.
	\problem{topic=quad}
	\problem{topic=ratineq}
	\problem{topic=graph}

	% An inline (one-off) problem can be placed right here:
	%
	% \begin{problem}{linear_eq}[topic=algebra, diff=easy]
	%     Solve $2x + 5 = 11$ for $x$.
	%     \solution
	%     $x = 3$.
	% \end{problem}

\end{problems}

% Optional: a final scoring grid for the front desk
% \scorepage[20]

\end{document}

%
% Build full collated multi-version PDF (no \Version):
%   lualatex autoexam-template.tex
%
% Build answer key:
%   lualatex \def\ShowKey{}% =============================================================================
% autoexam template — randomized, multi-version exam with key
% =============================================================================
%
% Build instructor copy (single version):
%   lualatex \def\Version{A}\input{autoexam-template.tex}
%
% Build full collated multi-version PDF (no \Version):
%   lualatex autoexam-template.tex
%
% Build answer key:
%   lualatex \def\ShowKey{}\input{autoexam-template.tex}
%
% Build with rubric overlays:
%   lualatex \def\ShowRubric{}\input{autoexam-template.tex}

\documentclass[
	exam-number     = 5,
	exam-date       = {May 2, 2026},
	exam-postscript = {Good luck!},
]{autoexam}

% -- Versions ------------------------------------------------------------------
% Declare the per-exam letter labels. The Python builder (if used) parses
% this line via regex to discover versions.
\versions{A, B, C}

% -- Problem bank --------------------------------------------------------------
% \loadbank pulls in problem definitions; alternatively, you can put
% \begin{problem}{id} ... \end{problem} blocks here (in the preamble or body).
\loadbank{bank.tex}

\begin{document}

% -- Cover page ----------------------------------------------------------------
% \maketitle emits the exam cover: the title (Exam <number>), course and date,
% the version circle (shown when \versions is declared), the boxed instructions
% (from autoexam-instructions.tex unless overridden), and the student name line.
% \begin{problems} issues a \clearpage after it, so Q1 starts on a fresh page.
% Add \blankpage here for a pre-problem scratch-work leaf.
\maketitle

% -- Body ----------------------------------------------------------------------
% Problems are grouped via \begin{problems} ... \end{problems}, which
% performs the smart two-column layout based on per-problem hints.
\begin{problems}

	% Inside \begin{problems}, use \problem{...} -- it emits the \question
	% wrapper that a problem's \begin{parts} needs. (\getproblem emits only
	% the bare problem body and is meant to follow an explicit \question,
	% e.g. in the quiz template.) The argument is a bank id or a key=value
	% filter; an optional [p1,p2,...] distributes points to \ppart calls.
	\problem{topic=quad}
	\problem{topic=ratineq}
	\problem{topic=graph}

	% An inline (one-off) problem can be placed right here:
	%
	% \begin{problem}{linear_eq}[topic=algebra, diff=easy]
	%     Solve $2x + 5 = 11$ for $x$.
	%     \solution
	%     $x = 3$.
	% \end{problem}

\end{problems}

% Optional: a final scoring grid for the front desk
% \scorepage[20]

\end{document}

%
% Build with rubric overlays:
%   lualatex \def\ShowRubric{}% =============================================================================
% autoexam template — randomized, multi-version exam with key
% =============================================================================
%
% Build instructor copy (single version):
%   lualatex \def\Version{A}\input{autoexam-template.tex}
%
% Build full collated multi-version PDF (no \Version):
%   lualatex autoexam-template.tex
%
% Build answer key:
%   lualatex \def\ShowKey{}\input{autoexam-template.tex}
%
% Build with rubric overlays:
%   lualatex \def\ShowRubric{}\input{autoexam-template.tex}

\documentclass[
	exam-number     = 5,
	exam-date       = {May 2, 2026},
	exam-postscript = {Good luck!},
]{autoexam}

% -- Versions ------------------------------------------------------------------
% Declare the per-exam letter labels. The Python builder (if used) parses
% this line via regex to discover versions.
\versions{A, B, C}

% -- Problem bank --------------------------------------------------------------
% \loadbank pulls in problem definitions; alternatively, you can put
% \begin{problem}{id} ... \end{problem} blocks here (in the preamble or body).
\loadbank{bank.tex}

\begin{document}

% -- Cover page ----------------------------------------------------------------
% \maketitle emits the exam cover: the title (Exam <number>), course and date,
% the version circle (shown when \versions is declared), the boxed instructions
% (from autoexam-instructions.tex unless overridden), and the student name line.
% \begin{problems} issues a \clearpage after it, so Q1 starts on a fresh page.
% Add \blankpage here for a pre-problem scratch-work leaf.
\maketitle

% -- Body ----------------------------------------------------------------------
% Problems are grouped via \begin{problems} ... \end{problems}, which
% performs the smart two-column layout based on per-problem hints.
\begin{problems}

	% Inside \begin{problems}, use \problem{...} -- it emits the \question
	% wrapper that a problem's \begin{parts} needs. (\getproblem emits only
	% the bare problem body and is meant to follow an explicit \question,
	% e.g. in the quiz template.) The argument is a bank id or a key=value
	% filter; an optional [p1,p2,...] distributes points to \ppart calls.
	\problem{topic=quad}
	\problem{topic=ratineq}
	\problem{topic=graph}

	% An inline (one-off) problem can be placed right here:
	%
	% \begin{problem}{linear_eq}[topic=algebra, diff=easy]
	%     Solve $2x + 5 = 11$ for $x$.
	%     \solution
	%     $x = 3$.
	% \end{problem}

\end{problems}

% Optional: a final scoring grid for the front desk
% \scorepage[20]

\end{document}


\documentclass[
	exam-number     = 5,
	exam-date       = {May 2, 2026},
	exam-postscript = {Good luck!},
]{autoexam}

% -- Versions ------------------------------------------------------------------
% Declare the per-exam letter labels. The Python builder (if used) parses
% this line via regex to discover versions.
\versions{A, B, C}

% -- Problem bank --------------------------------------------------------------
% \loadbank pulls in problem definitions; alternatively, you can put
% \begin{problem}{id} ... \end{problem} blocks here (in the preamble or body).
\loadbank{bank.tex}

\begin{document}

% -- Cover page ----------------------------------------------------------------
% \maketitle emits the exam cover: the title (Exam <number>), course and date,
% the version circle (shown when \versions is declared), the boxed instructions
% (from autoexam-instructions.tex unless overridden), and the student name line.
% \begin{problems} issues a \clearpage after it, so Q1 starts on a fresh page.
% Add \blankpage here for a pre-problem scratch-work leaf.
\maketitle

% -- Body ----------------------------------------------------------------------
% Problems are grouped via \begin{problems} ... \end{problems}, which
% performs the smart two-column layout based on per-problem hints.
\begin{problems}

	% Inside \begin{problems}, use \problem{...} -- it emits the \question
	% wrapper that a problem's \begin{parts} needs. (\getproblem emits only
	% the bare problem body and is meant to follow an explicit \question,
	% e.g. in the quiz template.) The argument is a bank id or a key=value
	% filter; an optional [p1,p2,...] distributes points to \ppart calls.
	\problem{topic=quad}
	\problem{topic=ratineq}
	\problem{topic=graph}

	% An inline (one-off) problem can be placed right here:
	%
	% \begin{problem}{linear_eq}[topic=algebra, diff=easy]
	%     Solve $2x + 5 = 11$ for $x$.
	%     \solution
	%     $x = 3$.
	% \end{problem}

\end{problems}

% Optional: a final scoring grid for the front desk
% \scorepage[20]

\end{document}

%
% Build full collated multi-version PDF (no \Version):
%   lualatex autoexam-template.tex
%
% Build answer key:
%   lualatex \def\ShowKey{}% =============================================================================
% autoexam template — randomized, multi-version exam with key
% =============================================================================
%
% Build instructor copy (single version):
%   lualatex \def\Version{A}% =============================================================================
% autoexam template — randomized, multi-version exam with key
% =============================================================================
%
% Build instructor copy (single version):
%   lualatex \def\Version{A}\input{autoexam-template.tex}
%
% Build full collated multi-version PDF (no \Version):
%   lualatex autoexam-template.tex
%
% Build answer key:
%   lualatex \def\ShowKey{}\input{autoexam-template.tex}
%
% Build with rubric overlays:
%   lualatex \def\ShowRubric{}\input{autoexam-template.tex}

\documentclass[
	exam-number     = 5,
	exam-date       = {May 2, 2026},
	exam-postscript = {Good luck!},
]{autoexam}

% -- Versions ------------------------------------------------------------------
% Declare the per-exam letter labels. The Python builder (if used) parses
% this line via regex to discover versions.
\versions{A, B, C}

% -- Problem bank --------------------------------------------------------------
% \loadbank pulls in problem definitions; alternatively, you can put
% \begin{problem}{id} ... \end{problem} blocks here (in the preamble or body).
\loadbank{bank.tex}

\begin{document}

% -- Cover page ----------------------------------------------------------------
% \maketitle emits the exam cover: the title (Exam <number>), course and date,
% the version circle (shown when \versions is declared), the boxed instructions
% (from autoexam-instructions.tex unless overridden), and the student name line.
% \begin{problems} issues a \clearpage after it, so Q1 starts on a fresh page.
% Add \blankpage here for a pre-problem scratch-work leaf.
\maketitle

% -- Body ----------------------------------------------------------------------
% Problems are grouped via \begin{problems} ... \end{problems}, which
% performs the smart two-column layout based on per-problem hints.
\begin{problems}

	% Inside \begin{problems}, use \problem{...} -- it emits the \question
	% wrapper that a problem's \begin{parts} needs. (\getproblem emits only
	% the bare problem body and is meant to follow an explicit \question,
	% e.g. in the quiz template.) The argument is a bank id or a key=value
	% filter; an optional [p1,p2,...] distributes points to \ppart calls.
	\problem{topic=quad}
	\problem{topic=ratineq}
	\problem{topic=graph}

	% An inline (one-off) problem can be placed right here:
	%
	% \begin{problem}{linear_eq}[topic=algebra, diff=easy]
	%     Solve $2x + 5 = 11$ for $x$.
	%     \solution
	%     $x = 3$.
	% \end{problem}

\end{problems}

% Optional: a final scoring grid for the front desk
% \scorepage[20]

\end{document}

%
% Build full collated multi-version PDF (no \Version):
%   lualatex autoexam-template.tex
%
% Build answer key:
%   lualatex \def\ShowKey{}% =============================================================================
% autoexam template — randomized, multi-version exam with key
% =============================================================================
%
% Build instructor copy (single version):
%   lualatex \def\Version{A}\input{autoexam-template.tex}
%
% Build full collated multi-version PDF (no \Version):
%   lualatex autoexam-template.tex
%
% Build answer key:
%   lualatex \def\ShowKey{}\input{autoexam-template.tex}
%
% Build with rubric overlays:
%   lualatex \def\ShowRubric{}\input{autoexam-template.tex}

\documentclass[
	exam-number     = 5,
	exam-date       = {May 2, 2026},
	exam-postscript = {Good luck!},
]{autoexam}

% -- Versions ------------------------------------------------------------------
% Declare the per-exam letter labels. The Python builder (if used) parses
% this line via regex to discover versions.
\versions{A, B, C}

% -- Problem bank --------------------------------------------------------------
% \loadbank pulls in problem definitions; alternatively, you can put
% \begin{problem}{id} ... \end{problem} blocks here (in the preamble or body).
\loadbank{bank.tex}

\begin{document}

% -- Cover page ----------------------------------------------------------------
% \maketitle emits the exam cover: the title (Exam <number>), course and date,
% the version circle (shown when \versions is declared), the boxed instructions
% (from autoexam-instructions.tex unless overridden), and the student name line.
% \begin{problems} issues a \clearpage after it, so Q1 starts on a fresh page.
% Add \blankpage here for a pre-problem scratch-work leaf.
\maketitle

% -- Body ----------------------------------------------------------------------
% Problems are grouped via \begin{problems} ... \end{problems}, which
% performs the smart two-column layout based on per-problem hints.
\begin{problems}

	% Inside \begin{problems}, use \problem{...} -- it emits the \question
	% wrapper that a problem's \begin{parts} needs. (\getproblem emits only
	% the bare problem body and is meant to follow an explicit \question,
	% e.g. in the quiz template.) The argument is a bank id or a key=value
	% filter; an optional [p1,p2,...] distributes points to \ppart calls.
	\problem{topic=quad}
	\problem{topic=ratineq}
	\problem{topic=graph}

	% An inline (one-off) problem can be placed right here:
	%
	% \begin{problem}{linear_eq}[topic=algebra, diff=easy]
	%     Solve $2x + 5 = 11$ for $x$.
	%     \solution
	%     $x = 3$.
	% \end{problem}

\end{problems}

% Optional: a final scoring grid for the front desk
% \scorepage[20]

\end{document}

%
% Build with rubric overlays:
%   lualatex \def\ShowRubric{}% =============================================================================
% autoexam template — randomized, multi-version exam with key
% =============================================================================
%
% Build instructor copy (single version):
%   lualatex \def\Version{A}\input{autoexam-template.tex}
%
% Build full collated multi-version PDF (no \Version):
%   lualatex autoexam-template.tex
%
% Build answer key:
%   lualatex \def\ShowKey{}\input{autoexam-template.tex}
%
% Build with rubric overlays:
%   lualatex \def\ShowRubric{}\input{autoexam-template.tex}

\documentclass[
	exam-number     = 5,
	exam-date       = {May 2, 2026},
	exam-postscript = {Good luck!},
]{autoexam}

% -- Versions ------------------------------------------------------------------
% Declare the per-exam letter labels. The Python builder (if used) parses
% this line via regex to discover versions.
\versions{A, B, C}

% -- Problem bank --------------------------------------------------------------
% \loadbank pulls in problem definitions; alternatively, you can put
% \begin{problem}{id} ... \end{problem} blocks here (in the preamble or body).
\loadbank{bank.tex}

\begin{document}

% -- Cover page ----------------------------------------------------------------
% \maketitle emits the exam cover: the title (Exam <number>), course and date,
% the version circle (shown when \versions is declared), the boxed instructions
% (from autoexam-instructions.tex unless overridden), and the student name line.
% \begin{problems} issues a \clearpage after it, so Q1 starts on a fresh page.
% Add \blankpage here for a pre-problem scratch-work leaf.
\maketitle

% -- Body ----------------------------------------------------------------------
% Problems are grouped via \begin{problems} ... \end{problems}, which
% performs the smart two-column layout based on per-problem hints.
\begin{problems}

	% Inside \begin{problems}, use \problem{...} -- it emits the \question
	% wrapper that a problem's \begin{parts} needs. (\getproblem emits only
	% the bare problem body and is meant to follow an explicit \question,
	% e.g. in the quiz template.) The argument is a bank id or a key=value
	% filter; an optional [p1,p2,...] distributes points to \ppart calls.
	\problem{topic=quad}
	\problem{topic=ratineq}
	\problem{topic=graph}

	% An inline (one-off) problem can be placed right here:
	%
	% \begin{problem}{linear_eq}[topic=algebra, diff=easy]
	%     Solve $2x + 5 = 11$ for $x$.
	%     \solution
	%     $x = 3$.
	% \end{problem}

\end{problems}

% Optional: a final scoring grid for the front desk
% \scorepage[20]

\end{document}


\documentclass[
	exam-number     = 5,
	exam-date       = {May 2, 2026},
	exam-postscript = {Good luck!},
]{autoexam}

% -- Versions ------------------------------------------------------------------
% Declare the per-exam letter labels. The Python builder (if used) parses
% this line via regex to discover versions.
\versions{A, B, C}

% -- Problem bank --------------------------------------------------------------
% \loadbank pulls in problem definitions; alternatively, you can put
% \begin{problem}{id} ... \end{problem} blocks here (in the preamble or body).
\loadbank{bank.tex}

\begin{document}

% -- Cover page ----------------------------------------------------------------
% \maketitle emits the exam cover: the title (Exam <number>), course and date,
% the version circle (shown when \versions is declared), the boxed instructions
% (from autoexam-instructions.tex unless overridden), and the student name line.
% \begin{problems} issues a \clearpage after it, so Q1 starts on a fresh page.
% Add \blankpage here for a pre-problem scratch-work leaf.
\maketitle

% -- Body ----------------------------------------------------------------------
% Problems are grouped via \begin{problems} ... \end{problems}, which
% performs the smart two-column layout based on per-problem hints.
\begin{problems}

	% Inside \begin{problems}, use \problem{...} -- it emits the \question
	% wrapper that a problem's \begin{parts} needs. (\getproblem emits only
	% the bare problem body and is meant to follow an explicit \question,
	% e.g. in the quiz template.) The argument is a bank id or a key=value
	% filter; an optional [p1,p2,...] distributes points to \ppart calls.
	\problem{topic=quad}
	\problem{topic=ratineq}
	\problem{topic=graph}

	% An inline (one-off) problem can be placed right here:
	%
	% \begin{problem}{linear_eq}[topic=algebra, diff=easy]
	%     Solve $2x + 5 = 11$ for $x$.
	%     \solution
	%     $x = 3$.
	% \end{problem}

\end{problems}

% Optional: a final scoring grid for the front desk
% \scorepage[20]

\end{document}

%
% Build with rubric overlays:
%   lualatex \def\ShowRubric{}% =============================================================================
% autoexam template — randomized, multi-version exam with key
% =============================================================================
%
% Build instructor copy (single version):
%   lualatex \def\Version{A}% =============================================================================
% autoexam template — randomized, multi-version exam with key
% =============================================================================
%
% Build instructor copy (single version):
%   lualatex \def\Version{A}\input{autoexam-template.tex}
%
% Build full collated multi-version PDF (no \Version):
%   lualatex autoexam-template.tex
%
% Build answer key:
%   lualatex \def\ShowKey{}\input{autoexam-template.tex}
%
% Build with rubric overlays:
%   lualatex \def\ShowRubric{}\input{autoexam-template.tex}

\documentclass[
	exam-number     = 5,
	exam-date       = {May 2, 2026},
	exam-postscript = {Good luck!},
]{autoexam}

% -- Versions ------------------------------------------------------------------
% Declare the per-exam letter labels. The Python builder (if used) parses
% this line via regex to discover versions.
\versions{A, B, C}

% -- Problem bank --------------------------------------------------------------
% \loadbank pulls in problem definitions; alternatively, you can put
% \begin{problem}{id} ... \end{problem} blocks here (in the preamble or body).
\loadbank{bank.tex}

\begin{document}

% -- Cover page ----------------------------------------------------------------
% \maketitle emits the exam cover: the title (Exam <number>), course and date,
% the version circle (shown when \versions is declared), the boxed instructions
% (from autoexam-instructions.tex unless overridden), and the student name line.
% \begin{problems} issues a \clearpage after it, so Q1 starts on a fresh page.
% Add \blankpage here for a pre-problem scratch-work leaf.
\maketitle

% -- Body ----------------------------------------------------------------------
% Problems are grouped via \begin{problems} ... \end{problems}, which
% performs the smart two-column layout based on per-problem hints.
\begin{problems}

	% Inside \begin{problems}, use \problem{...} -- it emits the \question
	% wrapper that a problem's \begin{parts} needs. (\getproblem emits only
	% the bare problem body and is meant to follow an explicit \question,
	% e.g. in the quiz template.) The argument is a bank id or a key=value
	% filter; an optional [p1,p2,...] distributes points to \ppart calls.
	\problem{topic=quad}
	\problem{topic=ratineq}
	\problem{topic=graph}

	% An inline (one-off) problem can be placed right here:
	%
	% \begin{problem}{linear_eq}[topic=algebra, diff=easy]
	%     Solve $2x + 5 = 11$ for $x$.
	%     \solution
	%     $x = 3$.
	% \end{problem}

\end{problems}

% Optional: a final scoring grid for the front desk
% \scorepage[20]

\end{document}

%
% Build full collated multi-version PDF (no \Version):
%   lualatex autoexam-template.tex
%
% Build answer key:
%   lualatex \def\ShowKey{}% =============================================================================
% autoexam template — randomized, multi-version exam with key
% =============================================================================
%
% Build instructor copy (single version):
%   lualatex \def\Version{A}\input{autoexam-template.tex}
%
% Build full collated multi-version PDF (no \Version):
%   lualatex autoexam-template.tex
%
% Build answer key:
%   lualatex \def\ShowKey{}\input{autoexam-template.tex}
%
% Build with rubric overlays:
%   lualatex \def\ShowRubric{}\input{autoexam-template.tex}

\documentclass[
	exam-number     = 5,
	exam-date       = {May 2, 2026},
	exam-postscript = {Good luck!},
]{autoexam}

% -- Versions ------------------------------------------------------------------
% Declare the per-exam letter labels. The Python builder (if used) parses
% this line via regex to discover versions.
\versions{A, B, C}

% -- Problem bank --------------------------------------------------------------
% \loadbank pulls in problem definitions; alternatively, you can put
% \begin{problem}{id} ... \end{problem} blocks here (in the preamble or body).
\loadbank{bank.tex}

\begin{document}

% -- Cover page ----------------------------------------------------------------
% \maketitle emits the exam cover: the title (Exam <number>), course and date,
% the version circle (shown when \versions is declared), the boxed instructions
% (from autoexam-instructions.tex unless overridden), and the student name line.
% \begin{problems} issues a \clearpage after it, so Q1 starts on a fresh page.
% Add \blankpage here for a pre-problem scratch-work leaf.
\maketitle

% -- Body ----------------------------------------------------------------------
% Problems are grouped via \begin{problems} ... \end{problems}, which
% performs the smart two-column layout based on per-problem hints.
\begin{problems}

	% Inside \begin{problems}, use \problem{...} -- it emits the \question
	% wrapper that a problem's \begin{parts} needs. (\getproblem emits only
	% the bare problem body and is meant to follow an explicit \question,
	% e.g. in the quiz template.) The argument is a bank id or a key=value
	% filter; an optional [p1,p2,...] distributes points to \ppart calls.
	\problem{topic=quad}
	\problem{topic=ratineq}
	\problem{topic=graph}

	% An inline (one-off) problem can be placed right here:
	%
	% \begin{problem}{linear_eq}[topic=algebra, diff=easy]
	%     Solve $2x + 5 = 11$ for $x$.
	%     \solution
	%     $x = 3$.
	% \end{problem}

\end{problems}

% Optional: a final scoring grid for the front desk
% \scorepage[20]

\end{document}

%
% Build with rubric overlays:
%   lualatex \def\ShowRubric{}% =============================================================================
% autoexam template — randomized, multi-version exam with key
% =============================================================================
%
% Build instructor copy (single version):
%   lualatex \def\Version{A}\input{autoexam-template.tex}
%
% Build full collated multi-version PDF (no \Version):
%   lualatex autoexam-template.tex
%
% Build answer key:
%   lualatex \def\ShowKey{}\input{autoexam-template.tex}
%
% Build with rubric overlays:
%   lualatex \def\ShowRubric{}\input{autoexam-template.tex}

\documentclass[
	exam-number     = 5,
	exam-date       = {May 2, 2026},
	exam-postscript = {Good luck!},
]{autoexam}

% -- Versions ------------------------------------------------------------------
% Declare the per-exam letter labels. The Python builder (if used) parses
% this line via regex to discover versions.
\versions{A, B, C}

% -- Problem bank --------------------------------------------------------------
% \loadbank pulls in problem definitions; alternatively, you can put
% \begin{problem}{id} ... \end{problem} blocks here (in the preamble or body).
\loadbank{bank.tex}

\begin{document}

% -- Cover page ----------------------------------------------------------------
% \maketitle emits the exam cover: the title (Exam <number>), course and date,
% the version circle (shown when \versions is declared), the boxed instructions
% (from autoexam-instructions.tex unless overridden), and the student name line.
% \begin{problems} issues a \clearpage after it, so Q1 starts on a fresh page.
% Add \blankpage here for a pre-problem scratch-work leaf.
\maketitle

% -- Body ----------------------------------------------------------------------
% Problems are grouped via \begin{problems} ... \end{problems}, which
% performs the smart two-column layout based on per-problem hints.
\begin{problems}

	% Inside \begin{problems}, use \problem{...} -- it emits the \question
	% wrapper that a problem's \begin{parts} needs. (\getproblem emits only
	% the bare problem body and is meant to follow an explicit \question,
	% e.g. in the quiz template.) The argument is a bank id or a key=value
	% filter; an optional [p1,p2,...] distributes points to \ppart calls.
	\problem{topic=quad}
	\problem{topic=ratineq}
	\problem{topic=graph}

	% An inline (one-off) problem can be placed right here:
	%
	% \begin{problem}{linear_eq}[topic=algebra, diff=easy]
	%     Solve $2x + 5 = 11$ for $x$.
	%     \solution
	%     $x = 3$.
	% \end{problem}

\end{problems}

% Optional: a final scoring grid for the front desk
% \scorepage[20]

\end{document}


\documentclass[
	exam-number     = 5,
	exam-date       = {May 2, 2026},
	exam-postscript = {Good luck!},
]{autoexam}

% -- Versions ------------------------------------------------------------------
% Declare the per-exam letter labels. The Python builder (if used) parses
% this line via regex to discover versions.
\versions{A, B, C}

% -- Problem bank --------------------------------------------------------------
% \loadbank pulls in problem definitions; alternatively, you can put
% \begin{problem}{id} ... \end{problem} blocks here (in the preamble or body).
\loadbank{bank.tex}

\begin{document}

% -- Cover page ----------------------------------------------------------------
% \maketitle emits the exam cover: the title (Exam <number>), course and date,
% the version circle (shown when \versions is declared), the boxed instructions
% (from autoexam-instructions.tex unless overridden), and the student name line.
% \begin{problems} issues a \clearpage after it, so Q1 starts on a fresh page.
% Add \blankpage here for a pre-problem scratch-work leaf.
\maketitle

% -- Body ----------------------------------------------------------------------
% Problems are grouped via \begin{problems} ... \end{problems}, which
% performs the smart two-column layout based on per-problem hints.
\begin{problems}

	% Inside \begin{problems}, use \problem{...} -- it emits the \question
	% wrapper that a problem's \begin{parts} needs. (\getproblem emits only
	% the bare problem body and is meant to follow an explicit \question,
	% e.g. in the quiz template.) The argument is a bank id or a key=value
	% filter; an optional [p1,p2,...] distributes points to \ppart calls.
	\problem{topic=quad}
	\problem{topic=ratineq}
	\problem{topic=graph}

	% An inline (one-off) problem can be placed right here:
	%
	% \begin{problem}{linear_eq}[topic=algebra, diff=easy]
	%     Solve $2x + 5 = 11$ for $x$.
	%     \solution
	%     $x = 3$.
	% \end{problem}

\end{problems}

% Optional: a final scoring grid for the front desk
% \scorepage[20]

\end{document}


\documentclass[
	exam-number     = 5,
	exam-date       = {May 2, 2026},
	exam-postscript = {Good luck!},
]{autoexam}

% -- Versions ------------------------------------------------------------------
% Declare the per-exam letter labels. The Python builder (if used) parses
% this line via regex to discover versions.
\versions{A, B, C}

% -- Problem bank --------------------------------------------------------------
% \loadbank pulls in problem definitions; alternatively, you can put
% \begin{problem}{id} ... \end{problem} blocks here (in the preamble or body).
\loadbank{bank.tex}

\begin{document}

% -- Cover page ----------------------------------------------------------------
% \maketitle emits the exam cover: the title (Exam <number>), course and date,
% the version circle (shown when \versions is declared), the boxed instructions
% (from autoexam-instructions.tex unless overridden), and the student name line.
% \begin{problems} issues a \clearpage after it, so Q1 starts on a fresh page.
% Add \blankpage here for a pre-problem scratch-work leaf.
\maketitle

% -- Body ----------------------------------------------------------------------
% Problems are grouped via \begin{problems} ... \end{problems}, which
% performs the smart two-column layout based on per-problem hints.
\begin{problems}

	% Inside \begin{problems}, use \problem{...} -- it emits the \question
	% wrapper that a problem's \begin{parts} needs. (\getproblem emits only
	% the bare problem body and is meant to follow an explicit \question,
	% e.g. in the quiz template.) The argument is a bank id or a key=value
	% filter; an optional [p1,p2,...] distributes points to \ppart calls.
	\problem{topic=quad}
	\problem{topic=ratineq}
	\problem{topic=graph}

	% An inline (one-off) problem can be placed right here:
	%
	% \begin{problem}{linear_eq}[topic=algebra, diff=easy]
	%     Solve $2x + 5 = 11$ for $x$.
	%     \solution
	%     $x = 3$.
	% \end{problem}

\end{problems}

% Optional: a final scoring grid for the front desk
% \scorepage[20]

\end{document}

%
% Build with rubric overlays:
%   lualatex \def\ShowRubric{}% =============================================================================
% autoexam template — randomized, multi-version exam with key
% =============================================================================
%
% Build instructor copy (single version):
%   lualatex \def\Version{A}% =============================================================================
% autoexam template — randomized, multi-version exam with key
% =============================================================================
%
% Build instructor copy (single version):
%   lualatex \def\Version{A}% =============================================================================
% autoexam template — randomized, multi-version exam with key
% =============================================================================
%
% Build instructor copy (single version):
%   lualatex \def\Version{A}\input{autoexam-template.tex}
%
% Build full collated multi-version PDF (no \Version):
%   lualatex autoexam-template.tex
%
% Build answer key:
%   lualatex \def\ShowKey{}\input{autoexam-template.tex}
%
% Build with rubric overlays:
%   lualatex \def\ShowRubric{}\input{autoexam-template.tex}

\documentclass[
	exam-number     = 5,
	exam-date       = {May 2, 2026},
	exam-postscript = {Good luck!},
]{autoexam}

% -- Versions ------------------------------------------------------------------
% Declare the per-exam letter labels. The Python builder (if used) parses
% this line via regex to discover versions.
\versions{A, B, C}

% -- Problem bank --------------------------------------------------------------
% \loadbank pulls in problem definitions; alternatively, you can put
% \begin{problem}{id} ... \end{problem} blocks here (in the preamble or body).
\loadbank{bank.tex}

\begin{document}

% -- Cover page ----------------------------------------------------------------
% \maketitle emits the exam cover: the title (Exam <number>), course and date,
% the version circle (shown when \versions is declared), the boxed instructions
% (from autoexam-instructions.tex unless overridden), and the student name line.
% \begin{problems} issues a \clearpage after it, so Q1 starts on a fresh page.
% Add \blankpage here for a pre-problem scratch-work leaf.
\maketitle

% -- Body ----------------------------------------------------------------------
% Problems are grouped via \begin{problems} ... \end{problems}, which
% performs the smart two-column layout based on per-problem hints.
\begin{problems}

	% Inside \begin{problems}, use \problem{...} -- it emits the \question
	% wrapper that a problem's \begin{parts} needs. (\getproblem emits only
	% the bare problem body and is meant to follow an explicit \question,
	% e.g. in the quiz template.) The argument is a bank id or a key=value
	% filter; an optional [p1,p2,...] distributes points to \ppart calls.
	\problem{topic=quad}
	\problem{topic=ratineq}
	\problem{topic=graph}

	% An inline (one-off) problem can be placed right here:
	%
	% \begin{problem}{linear_eq}[topic=algebra, diff=easy]
	%     Solve $2x + 5 = 11$ for $x$.
	%     \solution
	%     $x = 3$.
	% \end{problem}

\end{problems}

% Optional: a final scoring grid for the front desk
% \scorepage[20]

\end{document}

%
% Build full collated multi-version PDF (no \Version):
%   lualatex autoexam-template.tex
%
% Build answer key:
%   lualatex \def\ShowKey{}% =============================================================================
% autoexam template — randomized, multi-version exam with key
% =============================================================================
%
% Build instructor copy (single version):
%   lualatex \def\Version{A}\input{autoexam-template.tex}
%
% Build full collated multi-version PDF (no \Version):
%   lualatex autoexam-template.tex
%
% Build answer key:
%   lualatex \def\ShowKey{}\input{autoexam-template.tex}
%
% Build with rubric overlays:
%   lualatex \def\ShowRubric{}\input{autoexam-template.tex}

\documentclass[
	exam-number     = 5,
	exam-date       = {May 2, 2026},
	exam-postscript = {Good luck!},
]{autoexam}

% -- Versions ------------------------------------------------------------------
% Declare the per-exam letter labels. The Python builder (if used) parses
% this line via regex to discover versions.
\versions{A, B, C}

% -- Problem bank --------------------------------------------------------------
% \loadbank pulls in problem definitions; alternatively, you can put
% \begin{problem}{id} ... \end{problem} blocks here (in the preamble or body).
\loadbank{bank.tex}

\begin{document}

% -- Cover page ----------------------------------------------------------------
% \maketitle emits the exam cover: the title (Exam <number>), course and date,
% the version circle (shown when \versions is declared), the boxed instructions
% (from autoexam-instructions.tex unless overridden), and the student name line.
% \begin{problems} issues a \clearpage after it, so Q1 starts on a fresh page.
% Add \blankpage here for a pre-problem scratch-work leaf.
\maketitle

% -- Body ----------------------------------------------------------------------
% Problems are grouped via \begin{problems} ... \end{problems}, which
% performs the smart two-column layout based on per-problem hints.
\begin{problems}

	% Inside \begin{problems}, use \problem{...} -- it emits the \question
	% wrapper that a problem's \begin{parts} needs. (\getproblem emits only
	% the bare problem body and is meant to follow an explicit \question,
	% e.g. in the quiz template.) The argument is a bank id or a key=value
	% filter; an optional [p1,p2,...] distributes points to \ppart calls.
	\problem{topic=quad}
	\problem{topic=ratineq}
	\problem{topic=graph}

	% An inline (one-off) problem can be placed right here:
	%
	% \begin{problem}{linear_eq}[topic=algebra, diff=easy]
	%     Solve $2x + 5 = 11$ for $x$.
	%     \solution
	%     $x = 3$.
	% \end{problem}

\end{problems}

% Optional: a final scoring grid for the front desk
% \scorepage[20]

\end{document}

%
% Build with rubric overlays:
%   lualatex \def\ShowRubric{}% =============================================================================
% autoexam template — randomized, multi-version exam with key
% =============================================================================
%
% Build instructor copy (single version):
%   lualatex \def\Version{A}\input{autoexam-template.tex}
%
% Build full collated multi-version PDF (no \Version):
%   lualatex autoexam-template.tex
%
% Build answer key:
%   lualatex \def\ShowKey{}\input{autoexam-template.tex}
%
% Build with rubric overlays:
%   lualatex \def\ShowRubric{}\input{autoexam-template.tex}

\documentclass[
	exam-number     = 5,
	exam-date       = {May 2, 2026},
	exam-postscript = {Good luck!},
]{autoexam}

% -- Versions ------------------------------------------------------------------
% Declare the per-exam letter labels. The Python builder (if used) parses
% this line via regex to discover versions.
\versions{A, B, C}

% -- Problem bank --------------------------------------------------------------
% \loadbank pulls in problem definitions; alternatively, you can put
% \begin{problem}{id} ... \end{problem} blocks here (in the preamble or body).
\loadbank{bank.tex}

\begin{document}

% -- Cover page ----------------------------------------------------------------
% \maketitle emits the exam cover: the title (Exam <number>), course and date,
% the version circle (shown when \versions is declared), the boxed instructions
% (from autoexam-instructions.tex unless overridden), and the student name line.
% \begin{problems} issues a \clearpage after it, so Q1 starts on a fresh page.
% Add \blankpage here for a pre-problem scratch-work leaf.
\maketitle

% -- Body ----------------------------------------------------------------------
% Problems are grouped via \begin{problems} ... \end{problems}, which
% performs the smart two-column layout based on per-problem hints.
\begin{problems}

	% Inside \begin{problems}, use \problem{...} -- it emits the \question
	% wrapper that a problem's \begin{parts} needs. (\getproblem emits only
	% the bare problem body and is meant to follow an explicit \question,
	% e.g. in the quiz template.) The argument is a bank id or a key=value
	% filter; an optional [p1,p2,...] distributes points to \ppart calls.
	\problem{topic=quad}
	\problem{topic=ratineq}
	\problem{topic=graph}

	% An inline (one-off) problem can be placed right here:
	%
	% \begin{problem}{linear_eq}[topic=algebra, diff=easy]
	%     Solve $2x + 5 = 11$ for $x$.
	%     \solution
	%     $x = 3$.
	% \end{problem}

\end{problems}

% Optional: a final scoring grid for the front desk
% \scorepage[20]

\end{document}


\documentclass[
	exam-number     = 5,
	exam-date       = {May 2, 2026},
	exam-postscript = {Good luck!},
]{autoexam}

% -- Versions ------------------------------------------------------------------
% Declare the per-exam letter labels. The Python builder (if used) parses
% this line via regex to discover versions.
\versions{A, B, C}

% -- Problem bank --------------------------------------------------------------
% \loadbank pulls in problem definitions; alternatively, you can put
% \begin{problem}{id} ... \end{problem} blocks here (in the preamble or body).
\loadbank{bank.tex}

\begin{document}

% -- Cover page ----------------------------------------------------------------
% \maketitle emits the exam cover: the title (Exam <number>), course and date,
% the version circle (shown when \versions is declared), the boxed instructions
% (from autoexam-instructions.tex unless overridden), and the student name line.
% \begin{problems} issues a \clearpage after it, so Q1 starts on a fresh page.
% Add \blankpage here for a pre-problem scratch-work leaf.
\maketitle

% -- Body ----------------------------------------------------------------------
% Problems are grouped via \begin{problems} ... \end{problems}, which
% performs the smart two-column layout based on per-problem hints.
\begin{problems}

	% Inside \begin{problems}, use \problem{...} -- it emits the \question
	% wrapper that a problem's \begin{parts} needs. (\getproblem emits only
	% the bare problem body and is meant to follow an explicit \question,
	% e.g. in the quiz template.) The argument is a bank id or a key=value
	% filter; an optional [p1,p2,...] distributes points to \ppart calls.
	\problem{topic=quad}
	\problem{topic=ratineq}
	\problem{topic=graph}

	% An inline (one-off) problem can be placed right here:
	%
	% \begin{problem}{linear_eq}[topic=algebra, diff=easy]
	%     Solve $2x + 5 = 11$ for $x$.
	%     \solution
	%     $x = 3$.
	% \end{problem}

\end{problems}

% Optional: a final scoring grid for the front desk
% \scorepage[20]

\end{document}

%
% Build full collated multi-version PDF (no \Version):
%   lualatex autoexam-template.tex
%
% Build answer key:
%   lualatex \def\ShowKey{}% =============================================================================
% autoexam template — randomized, multi-version exam with key
% =============================================================================
%
% Build instructor copy (single version):
%   lualatex \def\Version{A}% =============================================================================
% autoexam template — randomized, multi-version exam with key
% =============================================================================
%
% Build instructor copy (single version):
%   lualatex \def\Version{A}\input{autoexam-template.tex}
%
% Build full collated multi-version PDF (no \Version):
%   lualatex autoexam-template.tex
%
% Build answer key:
%   lualatex \def\ShowKey{}\input{autoexam-template.tex}
%
% Build with rubric overlays:
%   lualatex \def\ShowRubric{}\input{autoexam-template.tex}

\documentclass[
	exam-number     = 5,
	exam-date       = {May 2, 2026},
	exam-postscript = {Good luck!},
]{autoexam}

% -- Versions ------------------------------------------------------------------
% Declare the per-exam letter labels. The Python builder (if used) parses
% this line via regex to discover versions.
\versions{A, B, C}

% -- Problem bank --------------------------------------------------------------
% \loadbank pulls in problem definitions; alternatively, you can put
% \begin{problem}{id} ... \end{problem} blocks here (in the preamble or body).
\loadbank{bank.tex}

\begin{document}

% -- Cover page ----------------------------------------------------------------
% \maketitle emits the exam cover: the title (Exam <number>), course and date,
% the version circle (shown when \versions is declared), the boxed instructions
% (from autoexam-instructions.tex unless overridden), and the student name line.
% \begin{problems} issues a \clearpage after it, so Q1 starts on a fresh page.
% Add \blankpage here for a pre-problem scratch-work leaf.
\maketitle

% -- Body ----------------------------------------------------------------------
% Problems are grouped via \begin{problems} ... \end{problems}, which
% performs the smart two-column layout based on per-problem hints.
\begin{problems}

	% Inside \begin{problems}, use \problem{...} -- it emits the \question
	% wrapper that a problem's \begin{parts} needs. (\getproblem emits only
	% the bare problem body and is meant to follow an explicit \question,
	% e.g. in the quiz template.) The argument is a bank id or a key=value
	% filter; an optional [p1,p2,...] distributes points to \ppart calls.
	\problem{topic=quad}
	\problem{topic=ratineq}
	\problem{topic=graph}

	% An inline (one-off) problem can be placed right here:
	%
	% \begin{problem}{linear_eq}[topic=algebra, diff=easy]
	%     Solve $2x + 5 = 11$ for $x$.
	%     \solution
	%     $x = 3$.
	% \end{problem}

\end{problems}

% Optional: a final scoring grid for the front desk
% \scorepage[20]

\end{document}

%
% Build full collated multi-version PDF (no \Version):
%   lualatex autoexam-template.tex
%
% Build answer key:
%   lualatex \def\ShowKey{}% =============================================================================
% autoexam template — randomized, multi-version exam with key
% =============================================================================
%
% Build instructor copy (single version):
%   lualatex \def\Version{A}\input{autoexam-template.tex}
%
% Build full collated multi-version PDF (no \Version):
%   lualatex autoexam-template.tex
%
% Build answer key:
%   lualatex \def\ShowKey{}\input{autoexam-template.tex}
%
% Build with rubric overlays:
%   lualatex \def\ShowRubric{}\input{autoexam-template.tex}

\documentclass[
	exam-number     = 5,
	exam-date       = {May 2, 2026},
	exam-postscript = {Good luck!},
]{autoexam}

% -- Versions ------------------------------------------------------------------
% Declare the per-exam letter labels. The Python builder (if used) parses
% this line via regex to discover versions.
\versions{A, B, C}

% -- Problem bank --------------------------------------------------------------
% \loadbank pulls in problem definitions; alternatively, you can put
% \begin{problem}{id} ... \end{problem} blocks here (in the preamble or body).
\loadbank{bank.tex}

\begin{document}

% -- Cover page ----------------------------------------------------------------
% \maketitle emits the exam cover: the title (Exam <number>), course and date,
% the version circle (shown when \versions is declared), the boxed instructions
% (from autoexam-instructions.tex unless overridden), and the student name line.
% \begin{problems} issues a \clearpage after it, so Q1 starts on a fresh page.
% Add \blankpage here for a pre-problem scratch-work leaf.
\maketitle

% -- Body ----------------------------------------------------------------------
% Problems are grouped via \begin{problems} ... \end{problems}, which
% performs the smart two-column layout based on per-problem hints.
\begin{problems}

	% Inside \begin{problems}, use \problem{...} -- it emits the \question
	% wrapper that a problem's \begin{parts} needs. (\getproblem emits only
	% the bare problem body and is meant to follow an explicit \question,
	% e.g. in the quiz template.) The argument is a bank id or a key=value
	% filter; an optional [p1,p2,...] distributes points to \ppart calls.
	\problem{topic=quad}
	\problem{topic=ratineq}
	\problem{topic=graph}

	% An inline (one-off) problem can be placed right here:
	%
	% \begin{problem}{linear_eq}[topic=algebra, diff=easy]
	%     Solve $2x + 5 = 11$ for $x$.
	%     \solution
	%     $x = 3$.
	% \end{problem}

\end{problems}

% Optional: a final scoring grid for the front desk
% \scorepage[20]

\end{document}

%
% Build with rubric overlays:
%   lualatex \def\ShowRubric{}% =============================================================================
% autoexam template — randomized, multi-version exam with key
% =============================================================================
%
% Build instructor copy (single version):
%   lualatex \def\Version{A}\input{autoexam-template.tex}
%
% Build full collated multi-version PDF (no \Version):
%   lualatex autoexam-template.tex
%
% Build answer key:
%   lualatex \def\ShowKey{}\input{autoexam-template.tex}
%
% Build with rubric overlays:
%   lualatex \def\ShowRubric{}\input{autoexam-template.tex}

\documentclass[
	exam-number     = 5,
	exam-date       = {May 2, 2026},
	exam-postscript = {Good luck!},
]{autoexam}

% -- Versions ------------------------------------------------------------------
% Declare the per-exam letter labels. The Python builder (if used) parses
% this line via regex to discover versions.
\versions{A, B, C}

% -- Problem bank --------------------------------------------------------------
% \loadbank pulls in problem definitions; alternatively, you can put
% \begin{problem}{id} ... \end{problem} blocks here (in the preamble or body).
\loadbank{bank.tex}

\begin{document}

% -- Cover page ----------------------------------------------------------------
% \maketitle emits the exam cover: the title (Exam <number>), course and date,
% the version circle (shown when \versions is declared), the boxed instructions
% (from autoexam-instructions.tex unless overridden), and the student name line.
% \begin{problems} issues a \clearpage after it, so Q1 starts on a fresh page.
% Add \blankpage here for a pre-problem scratch-work leaf.
\maketitle

% -- Body ----------------------------------------------------------------------
% Problems are grouped via \begin{problems} ... \end{problems}, which
% performs the smart two-column layout based on per-problem hints.
\begin{problems}

	% Inside \begin{problems}, use \problem{...} -- it emits the \question
	% wrapper that a problem's \begin{parts} needs. (\getproblem emits only
	% the bare problem body and is meant to follow an explicit \question,
	% e.g. in the quiz template.) The argument is a bank id or a key=value
	% filter; an optional [p1,p2,...] distributes points to \ppart calls.
	\problem{topic=quad}
	\problem{topic=ratineq}
	\problem{topic=graph}

	% An inline (one-off) problem can be placed right here:
	%
	% \begin{problem}{linear_eq}[topic=algebra, diff=easy]
	%     Solve $2x + 5 = 11$ for $x$.
	%     \solution
	%     $x = 3$.
	% \end{problem}

\end{problems}

% Optional: a final scoring grid for the front desk
% \scorepage[20]

\end{document}


\documentclass[
	exam-number     = 5,
	exam-date       = {May 2, 2026},
	exam-postscript = {Good luck!},
]{autoexam}

% -- Versions ------------------------------------------------------------------
% Declare the per-exam letter labels. The Python builder (if used) parses
% this line via regex to discover versions.
\versions{A, B, C}

% -- Problem bank --------------------------------------------------------------
% \loadbank pulls in problem definitions; alternatively, you can put
% \begin{problem}{id} ... \end{problem} blocks here (in the preamble or body).
\loadbank{bank.tex}

\begin{document}

% -- Cover page ----------------------------------------------------------------
% \maketitle emits the exam cover: the title (Exam <number>), course and date,
% the version circle (shown when \versions is declared), the boxed instructions
% (from autoexam-instructions.tex unless overridden), and the student name line.
% \begin{problems} issues a \clearpage after it, so Q1 starts on a fresh page.
% Add \blankpage here for a pre-problem scratch-work leaf.
\maketitle

% -- Body ----------------------------------------------------------------------
% Problems are grouped via \begin{problems} ... \end{problems}, which
% performs the smart two-column layout based on per-problem hints.
\begin{problems}

	% Inside \begin{problems}, use \problem{...} -- it emits the \question
	% wrapper that a problem's \begin{parts} needs. (\getproblem emits only
	% the bare problem body and is meant to follow an explicit \question,
	% e.g. in the quiz template.) The argument is a bank id or a key=value
	% filter; an optional [p1,p2,...] distributes points to \ppart calls.
	\problem{topic=quad}
	\problem{topic=ratineq}
	\problem{topic=graph}

	% An inline (one-off) problem can be placed right here:
	%
	% \begin{problem}{linear_eq}[topic=algebra, diff=easy]
	%     Solve $2x + 5 = 11$ for $x$.
	%     \solution
	%     $x = 3$.
	% \end{problem}

\end{problems}

% Optional: a final scoring grid for the front desk
% \scorepage[20]

\end{document}

%
% Build with rubric overlays:
%   lualatex \def\ShowRubric{}% =============================================================================
% autoexam template — randomized, multi-version exam with key
% =============================================================================
%
% Build instructor copy (single version):
%   lualatex \def\Version{A}% =============================================================================
% autoexam template — randomized, multi-version exam with key
% =============================================================================
%
% Build instructor copy (single version):
%   lualatex \def\Version{A}\input{autoexam-template.tex}
%
% Build full collated multi-version PDF (no \Version):
%   lualatex autoexam-template.tex
%
% Build answer key:
%   lualatex \def\ShowKey{}\input{autoexam-template.tex}
%
% Build with rubric overlays:
%   lualatex \def\ShowRubric{}\input{autoexam-template.tex}

\documentclass[
	exam-number     = 5,
	exam-date       = {May 2, 2026},
	exam-postscript = {Good luck!},
]{autoexam}

% -- Versions ------------------------------------------------------------------
% Declare the per-exam letter labels. The Python builder (if used) parses
% this line via regex to discover versions.
\versions{A, B, C}

% -- Problem bank --------------------------------------------------------------
% \loadbank pulls in problem definitions; alternatively, you can put
% \begin{problem}{id} ... \end{problem} blocks here (in the preamble or body).
\loadbank{bank.tex}

\begin{document}

% -- Cover page ----------------------------------------------------------------
% \maketitle emits the exam cover: the title (Exam <number>), course and date,
% the version circle (shown when \versions is declared), the boxed instructions
% (from autoexam-instructions.tex unless overridden), and the student name line.
% \begin{problems} issues a \clearpage after it, so Q1 starts on a fresh page.
% Add \blankpage here for a pre-problem scratch-work leaf.
\maketitle

% -- Body ----------------------------------------------------------------------
% Problems are grouped via \begin{problems} ... \end{problems}, which
% performs the smart two-column layout based on per-problem hints.
\begin{problems}

	% Inside \begin{problems}, use \problem{...} -- it emits the \question
	% wrapper that a problem's \begin{parts} needs. (\getproblem emits only
	% the bare problem body and is meant to follow an explicit \question,
	% e.g. in the quiz template.) The argument is a bank id or a key=value
	% filter; an optional [p1,p2,...] distributes points to \ppart calls.
	\problem{topic=quad}
	\problem{topic=ratineq}
	\problem{topic=graph}

	% An inline (one-off) problem can be placed right here:
	%
	% \begin{problem}{linear_eq}[topic=algebra, diff=easy]
	%     Solve $2x + 5 = 11$ for $x$.
	%     \solution
	%     $x = 3$.
	% \end{problem}

\end{problems}

% Optional: a final scoring grid for the front desk
% \scorepage[20]

\end{document}

%
% Build full collated multi-version PDF (no \Version):
%   lualatex autoexam-template.tex
%
% Build answer key:
%   lualatex \def\ShowKey{}% =============================================================================
% autoexam template — randomized, multi-version exam with key
% =============================================================================
%
% Build instructor copy (single version):
%   lualatex \def\Version{A}\input{autoexam-template.tex}
%
% Build full collated multi-version PDF (no \Version):
%   lualatex autoexam-template.tex
%
% Build answer key:
%   lualatex \def\ShowKey{}\input{autoexam-template.tex}
%
% Build with rubric overlays:
%   lualatex \def\ShowRubric{}\input{autoexam-template.tex}

\documentclass[
	exam-number     = 5,
	exam-date       = {May 2, 2026},
	exam-postscript = {Good luck!},
]{autoexam}

% -- Versions ------------------------------------------------------------------
% Declare the per-exam letter labels. The Python builder (if used) parses
% this line via regex to discover versions.
\versions{A, B, C}

% -- Problem bank --------------------------------------------------------------
% \loadbank pulls in problem definitions; alternatively, you can put
% \begin{problem}{id} ... \end{problem} blocks here (in the preamble or body).
\loadbank{bank.tex}

\begin{document}

% -- Cover page ----------------------------------------------------------------
% \maketitle emits the exam cover: the title (Exam <number>), course and date,
% the version circle (shown when \versions is declared), the boxed instructions
% (from autoexam-instructions.tex unless overridden), and the student name line.
% \begin{problems} issues a \clearpage after it, so Q1 starts on a fresh page.
% Add \blankpage here for a pre-problem scratch-work leaf.
\maketitle

% -- Body ----------------------------------------------------------------------
% Problems are grouped via \begin{problems} ... \end{problems}, which
% performs the smart two-column layout based on per-problem hints.
\begin{problems}

	% Inside \begin{problems}, use \problem{...} -- it emits the \question
	% wrapper that a problem's \begin{parts} needs. (\getproblem emits only
	% the bare problem body and is meant to follow an explicit \question,
	% e.g. in the quiz template.) The argument is a bank id or a key=value
	% filter; an optional [p1,p2,...] distributes points to \ppart calls.
	\problem{topic=quad}
	\problem{topic=ratineq}
	\problem{topic=graph}

	% An inline (one-off) problem can be placed right here:
	%
	% \begin{problem}{linear_eq}[topic=algebra, diff=easy]
	%     Solve $2x + 5 = 11$ for $x$.
	%     \solution
	%     $x = 3$.
	% \end{problem}

\end{problems}

% Optional: a final scoring grid for the front desk
% \scorepage[20]

\end{document}

%
% Build with rubric overlays:
%   lualatex \def\ShowRubric{}% =============================================================================
% autoexam template — randomized, multi-version exam with key
% =============================================================================
%
% Build instructor copy (single version):
%   lualatex \def\Version{A}\input{autoexam-template.tex}
%
% Build full collated multi-version PDF (no \Version):
%   lualatex autoexam-template.tex
%
% Build answer key:
%   lualatex \def\ShowKey{}\input{autoexam-template.tex}
%
% Build with rubric overlays:
%   lualatex \def\ShowRubric{}\input{autoexam-template.tex}

\documentclass[
	exam-number     = 5,
	exam-date       = {May 2, 2026},
	exam-postscript = {Good luck!},
]{autoexam}

% -- Versions ------------------------------------------------------------------
% Declare the per-exam letter labels. The Python builder (if used) parses
% this line via regex to discover versions.
\versions{A, B, C}

% -- Problem bank --------------------------------------------------------------
% \loadbank pulls in problem definitions; alternatively, you can put
% \begin{problem}{id} ... \end{problem} blocks here (in the preamble or body).
\loadbank{bank.tex}

\begin{document}

% -- Cover page ----------------------------------------------------------------
% \maketitle emits the exam cover: the title (Exam <number>), course and date,
% the version circle (shown when \versions is declared), the boxed instructions
% (from autoexam-instructions.tex unless overridden), and the student name line.
% \begin{problems} issues a \clearpage after it, so Q1 starts on a fresh page.
% Add \blankpage here for a pre-problem scratch-work leaf.
\maketitle

% -- Body ----------------------------------------------------------------------
% Problems are grouped via \begin{problems} ... \end{problems}, which
% performs the smart two-column layout based on per-problem hints.
\begin{problems}

	% Inside \begin{problems}, use \problem{...} -- it emits the \question
	% wrapper that a problem's \begin{parts} needs. (\getproblem emits only
	% the bare problem body and is meant to follow an explicit \question,
	% e.g. in the quiz template.) The argument is a bank id or a key=value
	% filter; an optional [p1,p2,...] distributes points to \ppart calls.
	\problem{topic=quad}
	\problem{topic=ratineq}
	\problem{topic=graph}

	% An inline (one-off) problem can be placed right here:
	%
	% \begin{problem}{linear_eq}[topic=algebra, diff=easy]
	%     Solve $2x + 5 = 11$ for $x$.
	%     \solution
	%     $x = 3$.
	% \end{problem}

\end{problems}

% Optional: a final scoring grid for the front desk
% \scorepage[20]

\end{document}


\documentclass[
	exam-number     = 5,
	exam-date       = {May 2, 2026},
	exam-postscript = {Good luck!},
]{autoexam}

% -- Versions ------------------------------------------------------------------
% Declare the per-exam letter labels. The Python builder (if used) parses
% this line via regex to discover versions.
\versions{A, B, C}

% -- Problem bank --------------------------------------------------------------
% \loadbank pulls in problem definitions; alternatively, you can put
% \begin{problem}{id} ... \end{problem} blocks here (in the preamble or body).
\loadbank{bank.tex}

\begin{document}

% -- Cover page ----------------------------------------------------------------
% \maketitle emits the exam cover: the title (Exam <number>), course and date,
% the version circle (shown when \versions is declared), the boxed instructions
% (from autoexam-instructions.tex unless overridden), and the student name line.
% \begin{problems} issues a \clearpage after it, so Q1 starts on a fresh page.
% Add \blankpage here for a pre-problem scratch-work leaf.
\maketitle

% -- Body ----------------------------------------------------------------------
% Problems are grouped via \begin{problems} ... \end{problems}, which
% performs the smart two-column layout based on per-problem hints.
\begin{problems}

	% Inside \begin{problems}, use \problem{...} -- it emits the \question
	% wrapper that a problem's \begin{parts} needs. (\getproblem emits only
	% the bare problem body and is meant to follow an explicit \question,
	% e.g. in the quiz template.) The argument is a bank id or a key=value
	% filter; an optional [p1,p2,...] distributes points to \ppart calls.
	\problem{topic=quad}
	\problem{topic=ratineq}
	\problem{topic=graph}

	% An inline (one-off) problem can be placed right here:
	%
	% \begin{problem}{linear_eq}[topic=algebra, diff=easy]
	%     Solve $2x + 5 = 11$ for $x$.
	%     \solution
	%     $x = 3$.
	% \end{problem}

\end{problems}

% Optional: a final scoring grid for the front desk
% \scorepage[20]

\end{document}


\documentclass[
	exam-number     = 5,
	exam-date       = {May 2, 2026},
	exam-postscript = {Good luck!},
]{autoexam}

% -- Versions ------------------------------------------------------------------
% Declare the per-exam letter labels. The Python builder (if used) parses
% this line via regex to discover versions.
\versions{A, B, C}

% -- Problem bank --------------------------------------------------------------
% \loadbank pulls in problem definitions; alternatively, you can put
% \begin{problem}{id} ... \end{problem} blocks here (in the preamble or body).
\loadbank{bank.tex}

\begin{document}

% -- Cover page ----------------------------------------------------------------
% \maketitle emits the exam cover: the title (Exam <number>), course and date,
% the version circle (shown when \versions is declared), the boxed instructions
% (from autoexam-instructions.tex unless overridden), and the student name line.
% \begin{problems} issues a \clearpage after it, so Q1 starts on a fresh page.
% Add \blankpage here for a pre-problem scratch-work leaf.
\maketitle

% -- Body ----------------------------------------------------------------------
% Problems are grouped via \begin{problems} ... \end{problems}, which
% performs the smart two-column layout based on per-problem hints.
\begin{problems}

	% Inside \begin{problems}, use \problem{...} -- it emits the \question
	% wrapper that a problem's \begin{parts} needs. (\getproblem emits only
	% the bare problem body and is meant to follow an explicit \question,
	% e.g. in the quiz template.) The argument is a bank id or a key=value
	% filter; an optional [p1,p2,...] distributes points to \ppart calls.
	\problem{topic=quad}
	\problem{topic=ratineq}
	\problem{topic=graph}

	% An inline (one-off) problem can be placed right here:
	%
	% \begin{problem}{linear_eq}[topic=algebra, diff=easy]
	%     Solve $2x + 5 = 11$ for $x$.
	%     \solution
	%     $x = 3$.
	% \end{problem}

\end{problems}

% Optional: a final scoring grid for the front desk
% \scorepage[20]

\end{document}


\documentclass[
	exam-number     = 5,
	exam-date       = {May 2, 2026},
	exam-postscript = {Good luck!},
]{autoexam}

% -- Versions ------------------------------------------------------------------
% Declare the per-exam letter labels. The Python builder (if used) parses
% this line via regex to discover versions.
\versions{A, B, C}

% -- Problem bank --------------------------------------------------------------
% \loadbank pulls in problem definitions; alternatively, you can put
% \begin{problem}{id} ... \end{problem} blocks here (in the preamble or body).
\loadbank{bank.tex}

\begin{document}

% -- Cover page ----------------------------------------------------------------
% \maketitle emits the exam cover: the title (Exam <number>), course and date,
% the version circle (shown when \versions is declared), the boxed instructions
% (from autoexam-instructions.tex unless overridden), and the student name line.
% \begin{problems} issues a \clearpage after it, so Q1 starts on a fresh page.
% Add \blankpage here for a pre-problem scratch-work leaf.
\maketitle

% -- Body ----------------------------------------------------------------------
% Problems are grouped via \begin{problems} ... \end{problems}, which
% performs the smart two-column layout based on per-problem hints.
\begin{problems}

	% Inside \begin{problems}, use \problem{...} -- it emits the \question
	% wrapper that a problem's \begin{parts} needs. (\getproblem emits only
	% the bare problem body and is meant to follow an explicit \question,
	% e.g. in the quiz template.) The argument is a bank id or a key=value
	% filter; an optional [p1,p2,...] distributes points to \ppart calls.
	\problem{topic=quad}
	\problem{topic=ratineq}
	\problem{topic=graph}

	% An inline (one-off) problem can be placed right here:
	%
	% \begin{problem}{linear_eq}[topic=algebra, diff=easy]
	%     Solve $2x + 5 = 11$ for $x$.
	%     \solution
	%     $x = 3$.
	% \end{problem}

\end{problems}

% Optional: a final scoring grid for the front desk
% \scorepage[20]

\end{document}
